\par\Needspace{6\baselineskip}\medskip
\hypertarget{native-083-article064}{}\pdfbookmark[2]{64.}{article-064}
\begin{center}64.\end{center}
We add the following observations on the problem of Article 62.

I. If in the second equation given there we set $N=\lambda$, $N=L$, $N=l$, we obtain $R'\sin(\lambda-L)=\Delta'\sin(l-\lambda)$, $r'\sin(\lambda-L)=\Delta'\sin(l-L)$, $r'\sin(l-\lambda)=R'\sin(l-L)$. The first or the second equation is convenient for checking the calculation when Method I or II of Article 62 has been used. In our example we thus have
\begin{center}
\begin{tabular}{@{}ll@{}}
\makebox[13em][l]{$\log\sin(\lambda-L)$\hfill .\,.\,.} & $9{,}475\,8653_n$\\[.3em]
\makebox[13em][l]{$\log\frac{\Delta'}{r'}$\hfill .\,.\,.\,.\,.\,.\,.\,.\,.} & $9{,}754\,6117$\\[.3em]
\hline
& $9{,}721\,2536_n$\\[.3em]
\makebox[13em][l]{$\log\sin(l-L)$\hfill .\,.\,.} & $9{,}721\,2536_n$
\end{tabular}
\end{center}
\begin{center}
$l-L=-31^0\,45^{\prime}\,26^{\prime\prime}{,}82$
\end{center}

II. The Sun and the two points in the plane of the ecliptic that are the projections of the positions of the celestial body and the Earth form a plane triangle with sides $\Delta'$, $R'$, $r'$ and opposite angles either $\lambda-L$, $l-\lambda$, $180^0-l+L$, or $L-\lambda$, $\lambda-l$, $180^0-L+l$. The relations given in I follow directly from this principle.

III. The Sun, the true position of the celestial body in space, and the true position of the Earth form another triangle, whose sides are $\Delta$, $R$, $r$. If we denote the respective opposite angles by $S$, $T$, $180^0-S-T$, we have $\frac{\sin S}{\Delta}=\frac{\sin T}{R}=\frac{\sin(S+T)}{r}\cdot$ The plane of this triangle traces a great circle on the celestial sphere. The heliocentric position of the Earth, the heliocentric position of the celestial body, and its geocentric position lie on this circle in such a way that the distances of the second from the first, the third from the second, and the third from the first, reckoned in the same direction, are respectively $S$, $T$, $S+T$.

\SourceLeaf{084}IV. The following differential equations follow either from the known differential variations of the parts of a plane triangle or just as easily from the formulas of Article 62.
\[
\mathrm{d}l=\frac{r'\cos(\lambda-l)}{\Delta'}\,\mathrm{d}\lambda+\frac{\sin(\lambda-l)}{\Delta'}\,\mathrm{d}r'
\]
\[
\mathrm{d}\Delta'=-r'\sin(\lambda-l)\,\mathrm{d}\lambda+\cos(\lambda-l)\,\mathrm{d}r'
\]
\[
\begin{aligned}
\mathrm{d}b={}&\frac{r'\cos b\sin b\sin(\lambda-l)}{\Delta'}\,\mathrm{d}\lambda
+\frac{r'\cos b^2}{\Delta'\cos\beta^2}\,\mathrm{d}\beta\\
&+\frac{\cos b^2}{\Delta'}(\operatorname{tang}\beta-\cos(\lambda-l)\operatorname{tang}b)\,\mathrm{d}r',
\end{aligned}
\]
\noindent Here, if the changes in the angles are expressed in arcseconds, the terms containing $\mathrm{d}r'$, $\mathrm{d}\Delta'$ must be multiplied by $206265^{\prime\prime}$, or the remaining terms divided by $206265^{\prime\prime}$.

V. The inverse problem of determining the heliocentric position from the geocentric one is entirely analogous to the problem developed above, so it would be superfluous to dwell on it further. All the formulas of Article 62 also apply to that problem, provided we interchange all the quantities relating to the heliocentric position of the celestial body with their counterparts for the geocentric position and substitute, for $L$, $B$ respectively, $L+180^0$, $-B$, or equivalently take the geocentric position of the Sun in place of the heliocentric position of the Earth.
